\subsection{値をポインタとして渡す；タグ付き共用体}

値をポインタで渡す方法の例です：

\begin{lstlisting}[label=unsigned_multiply_C,style=customc]
#include <stdio.h>
#include <stdint.h>

uint64_t multiply1 (uint64_t a, uint64_t b)
{
	return a*b;
};

uint64_t* multiply2 (uint64_t *a, uint64_t *b)
{
	return (uint64_t*)((uint64_t)a*(uint64_t)b);
};

int main()
{
	printf ("%d\n", multiply1(123, 456));
	printf ("%d\n", (uint64_t)multiply2((uint64_t*)123, (uint64_t*)456));
};
\end{lstlisting}

これはスムーズに動作し、GCC 4.8.4はmultiply1()とmultiply2()関数を同一にコンパイルします！

\begin{lstlisting}[label=unsigned_multiply_lst,style=customasmx86]
multiply1:
	mov	rax, rdi
	imul	rax, rsi
	ret

multiply2:
	mov	rax, rdi
	imul	rax, rsi
	ret
\end{lstlisting}

ポインタを逆参照しない限り（言い換えると、ポインタに格納されたアドレスからデータを読み取らない限り）、すべて正常に動作します。
ポインタは通常の変数のように何でも格納できる変数です。

\myindex{x86!\Instructions!MUL}
\myindex{x86!\Instructions!IMUL}
符号付き乗算命令（\IMUL）が符号なし乗算命令（\MUL）の代わりに使用されています。詳細については
\ref{IMUL_over_MUL}を参照してください。

\myindex{Tagged pointers}
ちなみに、ポインタを少し悪用する\IT{タグ付きポインタ}と呼ばれるよく知られたハックがあります。
簡単に言えば、すべてのポインタが例えば16バイトのサイズのメモリブロックを指している場合（または常に16バイト境界でアライメントされている場合）、ポインタの最下位4ビットは常にゼロビットであり、このスペースを何らかの方法で使用できます。
\myindex{LISP}
これはLISPコンパイラとインタープリターで非常に人気があります。
彼らはこれらの未使用ビットにセル/オブジェクトタイプを格納し、これによりメモリを節約できます。
さらに、追加のメモリアクセスなしに、ポインタだけでセル/オブジェクトタイプを判断できます。
詳細については：\InSqBrackets{\CNotes 1.3}。

% TODO Example of tagged ptrs here